% ============================================================================
\section{Marco teórico, conceptual y legal}

\subsection{Datos macroeconómicos en tiempo real}

Para evaluar con justicia qué tan bien predice un modelo, no se debe mirar el pasado usando cifras que recién se conocieron años después. En economía, cada boletín mensual ofrece una ``fotografía'' distinta de la historia. Si llamamos \(y_t^{(v)}\) al crecimiento publicado para el mes \(t\) dentro del boletín o \emph{vintage} \(v\), y llamamos \(\mathcal{I}_v\) al conjunto de toda la información ya divulgada hasta esa fecha, cualquier pronóstico hecho en ese momento solo puede usar datos dentro de \(\mathcal{I}_v\). Usar versiones corregidas posteriormente introduciría ventaja artificial por conocer el futuro.

\subsection{Revisiones estadísticas}

Una revisión estadística es simplemente el cambio que sufre un mismo dato económico entre dos boletines publicados en momentos distintos:
\begin{equation}
	r_{t}^{(v,v+h)}=y_t^{(v+h)}-y_t^{(v)}.
	\label{eq:revision}
\end{equation}
Cuando comparamos la publicación pasados \(h\) meses respecto a la \emph{primera vez} que ese dato salió a la luz (\(v_0(t)\)), hablamos de la revisión acumulada:
\begin{equation}
	R_{t,h}=y_t^{(v_0(t)+h)}-y_t^{(v_0(t))}.
	\label{eq:revision_acum}
\end{equation}

En todo el documento se distinguen tres ideas que suelen confundirse:
\begin{itemize}
	\item \textbf{Revisión:} cambio observado entre dos publicaciones de una misma cifra mensual del PBI.
	\item \textbf{Predicción de revisión:} anticipar cuánto cambiará una cifra ya anunciada, usando solo datos previos a su corrección.
	\item \textbf{Estimación del presente (\emph{nowcasting}):} estimar cómo creció la economía en el mes actual antes de que exista un anuncio oficial.
\end{itemize}
Este trabajo trata exclusivamente sobre la predicción de revisiones.

\subsection{Definición del objetivo y horizonte comparable}

Para comparar todas las fechas bajo una misma regla, la variable a predecir es la corrección registrada un número fijo de meses después del primer anuncio:
\begin{equation}
	Y_t^{(H)}=y_t^{(v_H(t))}-y_t^{(v_0(t))},
	\label{eq:target}
\end{equation}
donde \(v_0(t)\) es el primer boletín que informó el mes \(t\) y \(v_H(t)\) es el primer boletín publicado al menos \(H\) meses después. Aunque inicialmente se pensó usar doce meses (\(H=12\)), al revisar los archivos del BCRP se vio que la ventana móvil publicada históricamente deja \(18\) vacíos internos a ese plazo. Por ello se eligió \(H=11\) meses como medida anual continua principal (\(350\) observaciones sin huecos), conservando \(H=12\), \(6\) y \(3\) meses para verificar estabilidad.

No llamamos ``valor verdadero'' a la cifra revisada posterior, porque incluso años después pueden existir ajustes contables. Además, los \(24\) meses cruzados por cambios de año base se separaron del análisis principal para no confundir cambios de metodología con revisiones cotidianas.

\subsection{Modelos tradicionales de comparación}

Para saber si un método complejo realmente vale la pena, hay que compararlo con herramientas clásicas conocidas por su sencillez y solidez.

\subsubsection{Modelo autorregresivo de revisiones (AR y ARX)}

La idea más directa consiste en mirar si las revisiones anteriores ya completadas ayudan a predecir las nuevas:
\begin{equation}
	R_{t,H}=\alpha+\sum_{j=1}^{p}\phi_j R_{t-j,H}^{(obs)}+\boldsymbol{\beta}'\mathbf{x}_{t}^{(v_0)}+\varepsilon_t.
	\label{eq:ar}
\end{equation}
Aquí solo entran revisiones pasadas ya terminadas y variables ya publicadas en la fecha de pronóstico. El número de rezagos \(p\) se elige automáticamente dentro de cada historial pasado usando el criterio BIC, sin mirar nunca el tramo de prueba.

\subsubsection{Regresión regularizada (Ridge)}

Cuando se incluyen varias variables sectoriales a la vez, una regresión lineal común puede volverse inestable. La regresión Ridge \parencite{hoerl1970} evita coeficientes exagerados añadiendo un freno proporcional a su tamaño:
\begin{equation}
	\widehat{\boldsymbol{\beta}}=
	\arg\min_{\boldsymbol{\beta}}
	\left\{
	\sum_{t\in\mathcal{T}_{train}}(Y_t-\mathbf{x}_t'\boldsymbol{\beta})^2
	+\lambda\|\boldsymbol{\beta}\|_2^2
	\right\}.
	\label{eq:ridge}
\end{equation}
La intensidad del freno (\(\lambda\)) y la escala de las variables se calculan de nuevo en cada mes usando únicamente los datos pasados disponibles hasta ese instante.

\subsubsection{Árboles de decisión potenciados (GBDT)}

Los modelos lineales dibujan siempre rectas: si un dato cae mucho, proyectan una corrección gigante en proporción directa. En cambio, los árboles de decisión acotados \parencite{friedman2001,chen2016} dividen la información en reglas escalonadas tipo ``si ocurre tal rango, aplica este ajuste'', sumando varios árboles pequeños corregidos paso a paso. Esto les permite captar comportamientos curvos sin dispararse ante cifras extremas como las de una pandemia.

\subsubsection{Modelo de factores dinámicos (DFM)}

Este enfoque \parencite{stock2002,giannone2008} resume el pulso común de distintos sectores económicos en pocos factores principales:
\begin{align}
	\mathbf{x}_t&=\boldsymbol{\Lambda}f_t+\boldsymbol{\epsilon}_t,\\
	f_t&=\mathbf{A}f_{t-1}+\boldsymbol{\eta}_t.
\end{align}
Como la base histórica disponible reúne las tasas por sector económico (sin indicadores externos con fechas diarias certificadas), el modelo extrajo la tendencia común de las seis ramas productivas continuas y la proyectó mediante regresión Ridge.

\subsubsection{Corrección clásica por promedio histórico}

Como regla práctica elemental, se probó simplemente sumar a cada primer anuncio el promedio histórico de revisiones observadas hasta ese momento:
\begin{equation}
	\widehat{Y}_{t,H}^{\,BC}=\widehat{Y}_{t,H}^{\,base}+
	\widehat{\mathbb{E}}\left[R_{t,H}\mid\mathcal{I}_{v_0(t)}\right].
\end{equation}
Se evaluaron dos versiones: el promedio acumulado de toda la historia disponible y un promedio móvil de los últimos sesenta meses. Si un algoritmo complejo no logra ganar con claridad a esta regla tan simple, no justifica su uso.

\subsection{Computación de reservorio cuántico (QRC)}

Un reservorio funciona como un estanque con ondas: cuando ingresan los datos históricos \(u_t\), el sistema reacciona con una dinámica interna fija que mezcla la información de forma no lineal, sin necesidad de entrenar el interior del circuito \parencite{fujii2017}. Al final solo se toman \(m\) lecturas numéricas \(\mathbf{z}_t = (\langle O_1\rangle_t,\ldots,\langle O_m\rangle_t)'\) y se entrena una regresión Ridge sencilla para producir el pronóstico:
\begin{equation}
	\widehat{\mathbf{w}}=
	\arg\min_{\mathbf{w}}
	\left\{
	\sum_{t\in\mathcal{T}_{train}}
	(Y_t-\mathbf{w}'\mathbf{z}_t)^2+
	\lambda\|\mathbf{w}\|_2^2
	\right\}.
	\label{eq:qrc_readout}
\end{equation}
En este estudio se simuló de manera exacta (sin ruido físico de laboratorio) un sistema de cinco qubits que entrega veinte lecturas acotadas entre \(-1\) y \(+1\). Al dejar fija la parte cuántica y entrenar únicamente esa última línea recta sobre veinte señales, se evita el sobreajuste típico de circuitos cuánticos más profundos.

\subsection{Controles clásicos equivalentes}

Para comprobar si un buen resultado de QRC viene de ``lo cuántico'' o simplemente del tamaño y suavizado de sus veinte señales, se comparó en igualdad de condiciones contra dos métodos clásicos bien conocidos:
\begin{itemize}
	\item \textbf{Red de estado eco (ESN):} un reservorio clásico que también mezcla datos en el tiempo con pesos fijos y veinte neuronas de salida \parencite{jaeger2001} (más una versión mayor de sesenta y cuatro neuronas).
	\item \textbf{Características aleatorias de Fourier (RFF):} una transformación estática clásica que convierte los datos mediante veinte ondas aleatorias de seno y coseno \parencite{rahimi2007}, sin memoria recurrente interna.
\end{itemize}
Usar exactamente veinte señales de lectura garantiza una comparación honesta en complejidad y presupuesto de cálculo, sin asumir equivalencia matemática perfecta \parencite{schuld2021}.

\subsection{Qué no se puede afirmar sobre ventaja cuántica}

Incluso si QRC superara a una regresión lineal, eso no probaría ``supremacía cuántica'', porque una mejora podría deberse simplemente a que sus salidas están acotadas frente a una fórmula lineal que se dispara en crisis \parencite{schuld2021}. Por ello, este estudio evaluó estrictamente precisión real, estabilidad entre ejecuciones aleatorias y desempeño frente a controles clásicos equivalentes.

\subsection{Marco legal e institucional}

Todo el análisis empleó exclusivamente estadísticas públicas y código científico abierto, respetando licencias y atribución de origen. Las cuentas nacionales, los cambios de año base (2000 y 2014) y el calendario mensual de publicación se ajustan a las Notas Semanales del BCRP y al marco oficial del INEI para el PBI mensual.
